\chapter{ĐÓNG GÓP CỦA CÁC THÀNH VIÊN}
\label{chap:contrib}

% Bảng phân công dưới đây dựa trên cách chia 2 bài/người đã thống nhất.
% Nhóm điều chỉnh phần "Phần chung" và tỉ lệ đóng góp cho đúng thực tế.

Mỗi thành viên phụ trách đọc, phân tích và viết phần trình bày chi tiết cho hai nghiên cứu cùng hướng, đồng thời đảm nhận một phần chung của báo cáo. Bảng~\ref{tab:contrib} tóm tắt phân công.

\begin{table}[!htbp]
\centering
\small
\caption{Phân công công việc trong nhóm.}
\label{tab:contrib}
\begin{tabularx}{\textwidth}{L{3.0cm} L{3.6cm} Y C{1.4cm}}
\toprule
\textbf{Thành viên} & \textbf{Nghiên cứu phụ trách} & \textbf{Phần chung} & \textbf{Đóng góp} \\
\midrule
Trần Xuân Hương & \#1 MIRAGE, \#4 TIDE (Safety \& Robustness) & Tìm và chọn bài; Chương~1 (Giới thiệu); tổng hợp, biên tập báo cáo & 25\% \\
Vũ Tùng Dương & \#3 ProgRouter, \#5 HalluProp (Orchestration \& Control) & Chương~2 (Kiến thức nền); Chương~6 (Khoảng trống \& hướng phát triển) & 25\% \\
Phạm Việt Dũng & \#7 DAIN, \#8 AGENT-RADAR (Communication \& Interaction) & Chương~3 (Phân loại); Chương~5 (So sánh \& đánh giá) & 25\% \\
Lê Văn Đức & \#2 MAC-RRG, \#6 UrbanAgent (Domain Applications) & Chương~7 (Kết luận); thiết kế slide trình bày & 25\% \\
\bottomrule
\end{tabularx}
\end{table}

Tất cả thành viên cùng tham gia thảo luận khung phân loại, rà soát chéo số liệu giữa các phần và chuẩn bị thuyết trình.
